\subsection{刚有一点起色之后，如何防止再次滑落？}

如果你刚刚开始往正确方向移动，主要任务不是再用力一点去冲刺。
真正要做的，是在旧框架惯性仍然盘旋在旁边时，别让新的目标状态失去它刚刚得到的落脚点。
早期成果之所以容易塌掉，通常是因为下一次动作的障碍又被放回去了，而一个简单、可重复的常规还没有来得及装上。
保护成果的办法，是降低下一次动作的成本，稳定环境，并让下一次决策窗口比旧窗口更容易进入。

\begin{enumerate}[nosep]
  \item 明确你要保护的到底是什么早期成果。不要保护模糊意图，要保护一个可见变化，例如“我 8:30 开始学习了”或“这周下班后走路两次了”。
  \item 找出第一个会让旧模式回流的条件。通常是杂乱重新出现、模糊性回升、拖延、或疲劳，这些都会抬高障碍。
  \item 在下一次会话前先安装一个便宜稳定器。把书放桌上、摆好衣服、预先打开文件，或写好第一句，让目标状态从更强基础开始。
  \item 让下一次动作比上一次更小。早期成果脆弱，是因为意志能量仍然有限；更小的下一步能让系统继续前进，而不必付出昂贵重启成本。
  \item 至少连续几次重复同一条进入规则。重点是让新模式先显得普通，再给旧框架时间重新夺回场地。
  \item 提前检查下一次决策窗口。如果下一次开口点很清楚，就在它到来前做好准备，而不是等障碍重新升高后再即兴处理。
\end{enumerate}

\begin{remark}
假设你终于在下班后成功学了二十分钟。
成果不会靠“夸自己两句”就被保护住；它会靠“让明天的启动比今天更便宜”被保护住。
如果电脑保持打开、笔记保持可见、第一项任务已经写好，那么下一次会话就会以更低摩擦开始，新行为也就更有可能活下来。
\end{remark}